\chapter{2013年北京卷（文）}

\section{单选题}

\begin{problem}
已知集合 \(A = \{-1, 0, 1\}\)，\(B = \{x \mid -1 \leqslant x < 1\}\)，则 \(A \cap B =\)（\quad）

\choices
    {\(\{0\}\)}
    {\( \{-1,0\} \)}
    {\(\{0,1\}\)}
    {\(\{-1,0,1\}\)}
\end{problem}

\begin{answer}
B
\end{answer}

\begin{solution}
集合交集中的元素必须同时属于集合 \(A\) 和集合 \(B\). 在 \(A\) 的三个元素中，\(-1\) 和 \(0\) 满足 \(-1\leqslant x<1\)，而 \(1\) 不满足 \(x<1\)，所以 \(A\cap B=\{-1,0\}\)，应选 B.
\end{solution}

\begin{problem}
设 \(a\)，\(b\)，\(c \in \R\)，且 \(a > b\)，则（\quad）

\choices
    {\(ac > bc\)}
    {\(\frac{1}{a} \leqslant \frac{1}{b}\)}
    {\(a^2 > b^2\)}
    {\(a^3 > b^3\)}
\end{problem}

\begin{answer}
D
\end{answer}

\begin{solution}
逐项判断. 取 \(a=1\)，\(b=0\)，\(c=-1\)，则 \(a>b\)，但 \(ac=-1<0=bc\)，所以 A 错. 取 \(a=1\)，\(b=-1\)，则 \(a>b\)，但 \(\frac{1}{a}=1>-1=\frac{1}{b}\)，所以 B 错. 取 \(a=-1\)，\(b=-2\)，则 \(a>b\)，但 \(a^2=1<4=b^2\)，所以 C 错.

又因为 \(a>b\)，且函数 \(y=x^3\) 在 \(\R\) 上单调递增，所以 \(a^3>b^3\). 也可写成 \(a^3-b^3=(a-b)(a^2+ab+b^2)>0\)，故应选 D.
\end{solution}

\begin{problem}
下列函数中，既是偶函数又在区间 \((0, +\infty)\) 上单调递减的是（\quad）

\choices
    {\(y = \frac{1}{x}\)}
    {\( y = \e^{-x} \)}
    {\( y = -x^{2} + 1 \)}
    {\(y = \lg |x|\)}
\end{problem}

\begin{answer}
C
\end{answer}

\begin{solution}
选项 A 的函数 \(y=\frac{1}{x}\) 满足 \(f(-x)=-f(x)\)，是奇函数，不是偶函数. 选项 B 中 \(f(-x)=\e^x\neq\e^{-x}=f(x)\)，不是偶函数. 选项 C 中 \(f(-x)=-(-x)^2+1=f(x)\)，是偶函数；当 \(0<x_1<x_2\) 时，\(-x_1^2+1>-x_2^2+1\)，所以它在 \((0,+\infty)\) 上单调递减. 选项 D 虽为偶函数，但 \(\lg x\) 在 \((0,1)\) 上递减、在 \((1,+\infty)\) 上递增，不满足题意. 因此应选 C.
\end{solution}

\begin{problem}
在复平面内，复数 \( \i(2 - \i) \) 对应的点位于（\quad）

\choices
    {第一象限}
    {第二象限}
    {第三象限}
    {第四象限}
\end{problem}

\begin{answer}
A
\end{answer}

\begin{solution}
计算得

\(\i(2-\i)=2\i-\i^2=1+2\i\).

因此该复数在复平面内对应的点为 \((1,2)\)，实部和虚部都为正，点位于第一象限，应选 A.
\end{solution}

\begin{problem}
在 \(\triangle ABC\) 中，\(a = 3\)，\(b = 5\)，\(\sin A = \frac{1}{3}\)，则 \(\sin B =\)（\quad）

\choices
    {\( \frac{1}{5} \)}
    {\( \frac{5}{9} \)}
    {\( \frac{\sqrt{5}}{3} \)}
    {\(1\)}
\end{problem}

\begin{answer}
B
\end{answer}

\begin{solution}
由正弦定理，

\(\frac{a}{\sin A}=\frac{b}{\sin B}\).

代入 \(a=3\)，\(b=5\)，\(\sin A=\frac{1}{3}\)，得

\(\sin B=\frac{b}{a}\sin A=\frac{5}{3}\cdot\frac{1}{3}=\frac{5}{9}\).

故应选 B.
\end{solution}

\begin{problem}
执行如图所示的程序框图，输出的 \(S\) 值为（\quad）

\texfigure[max width=0.30\linewidth]{img_repaint/2013/beijing_liberal/q6_fig1.tex}

\choices
    {\(1\)}
    {\( \frac{2}{3} \)}
    {\( \frac{13}{21} \)}
    {\( \frac{610}{987} \)}
\end{problem}

\begin{answer}
C
\end{answer}

\begin{solution}
初始 \(i=0\)，\(S=1\). 第一次循环先计算

\(S=\frac{1^2+1}{2\cdot1+1}=\frac{2}{3}\)，再令 \(i=1\). 此时 \(i<2\)，继续循环.

第二次循环计算

\(S=\frac{(\frac{2}{3})^2+1}{2\cdot\frac{2}{3}+1}=\frac{\frac{4}{9}+1}{\frac{4}{3}+1}=\frac{13}{21}\)，再令 \(i=2\). 此时 \(i\geqslant2\)，循环结束，输出 \(S=\frac{13}{21}\)，应选 C.
\end{solution}

\begin{problem}
双曲线 \(x^{2} - \frac{y^{2}}{m} = 1\) 的离心率大于 \(\sqrt{2}\) 的充分必要条件是（\quad）

\choices
    {\(m > \frac{1}{2}\)}
    {\(m \geqslant 1\)}
    {\(m > 1\)}
    {\(m > 2\)}
\end{problem}

\begin{answer}
C
\end{answer}

\begin{solution}
由于题目给出的是双曲线，必有 \(m>0\). 将方程写成标准形式

\(\frac{x^2}{1^2}-\frac{y^2}{(\sqrt{m})^2}=1\).

因此半实轴 \(a=1\)，半虚轴 \(b=\sqrt{m}\)，焦半距满足 \(c^2=a^2+b^2=1+m\)，离心率为 \(e=\frac{c}{a}=\sqrt{1+m}\). 于是

\(e>\sqrt{2}\iff 1+m>2\iff m>1\).

所以离心率大于 \(\sqrt{2}\) 的充分必要条件是 \(m>1\)，应选 C.
\end{solution}

\begin{problem}
如图，在正方体 \(ABCD - A_{1}B_{1}C_{1}D_{1}\) 中，\(P\) 为对角线 \(BD_{1}\) 的三等分点，\(P\) 到各顶点的距离的不同取值有（\quad）

\texfigure[max width=0.32\linewidth]{img_repaint/2013/beijing_liberal/q8_fig1.tex}

\choices
    {\(3\) 个}
    {\(4\) 个}
    {\(5\) 个}
    {\(6\) 个}
\end{problem}

\begin{answer}
B
\end{answer}

\begin{solution}
设正方体棱长为 \(a\)，取坐标
\(D=(0,0,0)\)，\(A=(a,0,0)\)，\(B=(a,a,0)\)，\(C=(0,a,0)\)，\(D_1=(0,0,a)\)，并取 \(P\) 为从 \(B\) 到 \(D_1\) 的三等分点之一，则

\(P=B+\frac{1}{3}(D_1-B)=\left(\frac{2a}{3},\frac{2a}{3},\frac{a}{3}\right)\).

由两点间距离公式，

\(PA^2=PC^2=PB_1^2=\frac{2a^2}{3}\)，
\(PD^2=PA_1^2=PC_1^2=a^2\)，
\(PB^2=\frac{a^2}{3}\)，
\(PD_1^2=\frac{4a^2}{3}\).

所以不同的距离为 \(\frac{a}{\sqrt{3}}\)、\(a\sqrt{\frac{2}{3}}\)、\(a\)、\(\frac{2a}{\sqrt{3}}\)，共 \(4\) 个. 另一三等分点与此情形对称，所得不同距离个数相同，故应选 B.
\end{solution}

\section{填空题}

\begin{problem}
若抛物线 \( y^{2}=2px \) 的焦点坐标为 \( (1,0) \)，则 \(p=\)\fillinblank{}；准线方程为\fillinblank{}.
\end{problem}

\begin{answer}
\(2\)，\(x=-1\)
\end{answer}

\begin{solution}
抛物线 \(y^2=2px\) 的焦点为 \(\left(\frac{p}{2},0\right)\)，由焦点坐标为 \((1,0)\) 得 \(\frac{p}{2}=1\)，所以 \(p=2\). 此时抛物线为 \(y^2=4x\)，其准线方程为 \(x=-1\).
\end{solution}

\begin{problem}
某四棱锥的三视图如图所示，该四棱锥的体积为\fillinblank{}.

\texfigure[max width=0.34\linewidth]{img_repaint/2013/beijing_liberal/q10_fig1.tex}
\end{problem}

\begin{answer}
\(3\)
\end{answer}

\begin{solution}
由正视图可知，底面在该方向的长度为 \(2+1=3\)；由侧视图可知，底面在另一方向的长度为 \(1+2=3\). 结合俯视图可知底面是边长为 \(3\) 的正方形，且正视图、侧视图的高均为 \(1\)，所以四棱锥的高为 \(1\). 因而底面积为 \(3^2=9\)，四棱锥体积为 \(V=\frac{1}{3}\times9\times1=3\).
\end{solution}

\begin{problem}
若等比数列 \(\{a_{n}\}\) 满足 \(a_{2} + a_{4} = 20\)，\(a_{3} + a_{5} = 40\)，则公比 \(q = \)\fillinblank{} . 前 \(n\) 项和 \(S_{n} =\)\fillinblank{}.
\end{problem}

\begin{answer}
\(2\)，\(2^{n+1}-2\)
\end{answer}

\begin{solution}
设等比数列的公比为 \(q\). 由 \(a_4=a_2q^2\)，\(a_5=a_3q^2\)，得 \(a_2(1+q^2)=20\)，\(a_3(1+q^2)=40\). 由于 \(1+q^2>0\)，且第一式右端非零，\(a_2\neq0\)，两式相除得 \(\frac{a_3}{a_2}=2\)，即 \(q=2\).

代回第一式，得 \(a_2(1+4)=20\)，所以 \(a_2=4\)，从而 \(a_1=\frac{a_2}{q}=2\). 因此 \(S_n=\frac{a_1(q^n-1)}{q-1}=\frac{2(2^n-1)}{1}=2^{n+1}-2\). 故 \(q=2\)，\(S_n=2^{n+1}-2\).
\end{solution}

\begin{problem}
设 \(D\) 为不等式组 \(\begin{cases}
x \geqslant 0,\\
2x - y \leqslant 0,\\
x + y - 3 \leqslant 0
\end{cases}\) 表示的平面区域. 区域 \(D\) 上的点与点 \((1,0)\) 之间的距离的最小值为\fillinblank{}.
\end{problem}

\begin{answer}
\(\frac{2\sqrt{5}}{5}\)
\end{answer}

\begin{solution}
不等式组给出的区域满足 \(2x-y\leqslant0\)，即位于直线 \(2x-y=0\) 的一侧. 点 \((1,0)\) 到该直线的距离为 \(d=\frac{|2\cdot1-0|}{\sqrt{2^2+(-1)^2}}=\frac{2\sqrt{5}}{5}\).

从点 \((1,0)\) 向直线 \(2x-y=0\) 作垂线，垂足为 \(\left(1,0\right)-\frac{2}{5}(2,-1)=\left(\frac{1}{5},\frac{2}{5}\right)\). 该点满足 \(x\geqslant0\)、\(2x-y=0\) 以及 \(x+y-3=-\frac{12}{5}\leqslant0\)，所以垂足属于区域 \(D\). 因此区域内的最小距离正好等于上述点到直线的距离，即 \(\frac{2\sqrt{5}}{5}\).
\end{solution}

\begin{problem}
函数 \( f(x)=\begin{cases}
\log_{\frac{1}{2}}x,&x\geqslant1,\\
2^{x},&x<1
\end{cases} \) 的值域为\fillinblank{}.
\end{problem}

\begin{answer}
\((-\infty,2)\)
\end{answer}

\begin{solution}
当 \(x\geqslant1\) 时，\(f(x)=\log_{\frac{1}{2}}x\). 因为底数 \(\frac{1}{2}\) 介于 \(0\) 和 \(1\) 之间，且 \(x\geqslant1\)，所以该部分的值域为 \((-\infty,0]\).

当 \(x<1\) 时，\(f(x)=2^x\). 指数函数单调递增，且 \(x<1\)，所以该部分的值域为 \((0,2)\). 两部分值域的并集为 \((-\infty,2)\).
\end{solution}

\begin{problem}
已知点 \(A(1, -1)\)，\(B(3, 0)\)，\(C(2, 1)\). 若平面区域 \(D\) 由所有满足 \(\overrightarrow{AP} = \lambda \overrightarrow{AB} + \mu \overrightarrow{AC}\) （\(1 \leqslant \lambda \leqslant 2\)，\(0 \leqslant \mu \leqslant 1\)） 的点 \(P\) 组成. 则 \(D\) 的面积为\fillinblank{}.
\end{problem}

\begin{answer}
\(3\)
\end{answer}

\begin{solution}
由点的坐标得 \(\overrightarrow{AB}=(2,1)\)，\(\overrightarrow{AC}=(1,2)\). 任意点 \(P\) 可表示为 \(P=A+\lambda\overrightarrow{AB}+\mu\overrightarrow{AC}\)，其中 \(\lambda\) 的取值区间长度为 \(1\)，\(\mu\) 的取值区间长度也为 \(1\). 因此区域 \(D\) 是由两个边向量 \(\overrightarrow{AB}\)、\(\overrightarrow{AC}\) 张成的平行四边形，其面积为

\(S=\left|\det\begin{pmatrix}2&1\\1&2\end{pmatrix}\right|=|4-1|=3\).

故区域 \(D\) 的面积为 \(3\).
\end{solution}

\section{解答题}

\begin{problem}
已知函数 \( f(x) = (2\cos^2 x - 1)\sin 2x + \frac{1}{2}\cos 4x \).
\begin{enumerate}
\item 求 \( f(x) \) 的最小正周期及最大值；
\item 若 \(\alpha \in \left(\frac{\pi}{2}, \pi\right)\)，且 \(f(\alpha) = \frac{\sqrt{2}}{2}\)，求 \(\alpha\) 的值.
\end{enumerate}
\end{problem}

\begin{solution}
\begin{enumerate}
\item 由 \(2\cos^2 x-1=\cos 2x\)，得
\(f(x)=\cos 2x\sin 2x+\frac{1}{2}\cos 4x=\frac{1}{2}\sin 4x+\frac{1}{2}\cos 4x=\frac{\sqrt{2}}{2}\sin\left(4x+\frac{\pi}{4}\right)\).
因此，\(f(x)\) 的最小正周期为 \(T=\frac{2\pi}{4}=\frac{\pi}{2}\)，最大值为 \(\frac{\sqrt{2}}{2}\).
\item 由 \(f(\alpha)=\frac{\sqrt{2}}{2}\)，得 \(\sin\left(4\alpha+\frac{\pi}{4}\right)=1\)，所以 \(4\alpha+\frac{\pi}{4}=\frac{\pi}{2}+2k\pi\)，其中 \(k\in\Z\). 从而 \(\alpha=\frac{\pi}{16}+\frac{k\pi}{2}\). 因为 \(\frac{\pi}{2}<\alpha<\pi\)，所以 \(7<8k<15\)，故 \(k=1\)，即 \(\alpha=\frac{9\pi}{16}\).
\end{enumerate}
\end{solution}

\begin{problem}
下图是某市3月1日至14日的空气质量指数趋势图，空气质量指数小于100表示空气质量优良，空气质量指数大于200表示空气重度污染. 某人随机选择3月1日至3月13日中的某一天到达该市，并停留2天.

\begin{enumerate}
\item 求此人到达当日空气质量优良的概率；
\item 求此人在该市停留期间只有 1 天空气质量重度污染的概率；
\item 由图判断从哪天开始连续三天的空气质量指数方差最大？（结论不要求证明）
\end{enumerate}

\texfigure[max width=0.62\linewidth]{img_repaint/2013/beijing_liberal/q16_fig1.tex}
\end{problem}

\begin{solution}
\begin{enumerate}
\item 由图可读得3月1日至3月13日的空气质量指数依次为 \(86\)，\(25\)，\(57\)，\(143\)，\(220\)，\(160\)，\(40\)，\(217\)，\(160\)，\(121\)，\(158\)，\(86\)，\(79\). 其中小于 \(100\) 的日期组成的集合为 \(\{1,2,3,7,12,13\}\)，共有 \(6\) 天. 因为到达日等可能地取3月1日至3月13日中的一天，所以所求概率为 \(\frac{6}{13}\).
\item 由图可知空气质量指数大于 \(200\) 的日期为5日和8日. 到达日分别为1日至13日时，停留期间的指数对依次为 \((86,25)\)，\((25,57)\)，\((57,143)\)，\((143,220)\)，\((220,160)\)，\((160,40)\)，\((40,217)\)，\((217,160)\)，\((160,121)\)，\((121,158)\)，\((158,86)\)，\((86,79)\)，\((79,37)\). 其中恰有一个指数大于 \(200\) 的情形有4种，即到达日为4日、5日、7日或8日. 因此所求概率为 \(\frac{4}{13}\).
\item 对每组三天的数据计算方差并比较可知，从5日开始的连续三天数据 \(220\)，\(160\)，\(40\) 的方差最大. 例如该组三天的平均数为 \(140\)，方差为 \(\frac{(220-140)^2+(160-140)^2+(40-140)^2}{3}=5600\)，故结论为从5日开始.
\end{enumerate}
\end{solution}

\begin{problem}
如图，在四棱锥 \(P - ABCD\) 中，\(AB \parallel CD\)，\(AB \perp AD\)，\(CD = 2AB\)，平面 \(PAD \perp\) 底面 \(ABCD\)，\(PA \perp AD\)，\(E\) 和 \(F\) 分别是 \(CD\) 和 \(PC\) 的中点.

求证：
\begin{enumerate}
\item \(PA \perp\) 底面 \(ABCD\)；
\item \(BE \parallel\) 平面 \(PAD\)；
\item 平面 \(BEF \perp\) 平面 \(PCD\).
\end{enumerate}

\texfigure[max width=0.34\linewidth]{img_repaint/2013/beijing_liberal/q17_fig1.tex}
\end{problem}

\begin{solution}
\begin{enumerate}
\item 平面 \(PAD\) 与平面 \(ABCD\) 垂直，且两平面的交线为 \(AD\). 因为 \(PA\) 在平面 \(PAD\) 内且 \(PA\perp AD\)，由两平面垂直的性质定理，得 \(PA\perp\) 底面 \(ABCD\).
\item 因为 \(E\) 是 \(CD\) 的中点，所以 \(DE=\frac{1}{2}CD=AB\). 又因为 \(DE\parallel CD\parallel AB\)，所以四边形 \(ABED\) 是平行四边形，从而 \(BE\parallel AD\). 平面 \(PAD\) 与底面 \(ABCD\) 的交线为 \(AD\)，而 \(B\) 在底面内且不在 \(AD\) 上，所以 \(BE\) 不在平面 \(PAD\) 内. 由于 \(AD\subset\) 平面 \(PAD\)，故 \(BE\parallel\) 平面 \(PAD\).
\item 由 \(BE\parallel AD\) 及 \(AB\perp AD\)，得 \(BE\perp AB\). 又 \(CD\parallel AB\)，所以 \(BE\perp CD\). 由第（1）问知 \(PA\perp\) 底面 \(ABCD\)，故 \(PA\perp AB\). 于是 \(AB\) 垂直于平面 \(PAD\) 内相交的两条直线 \(PA\) 和 \(AD\)，从而 \(AB\perp\) 平面 \(PAD\). 因为 \(CD\parallel AB\)，所以 \(CD\perp\) 平面 \(PAD\)，进而 \(CD\perp PD\).

在三角形 \(PCD\) 中，\(E\)，\(F\) 分别是 \(CD\)，\(PC\) 的中点，所以 \(EF\parallel PD\)，从而 \(CD\perp EF\). 直线 \(BE\) 与 \(EF\) 是平面 \(BEF\) 内相交的两条直线，且都垂直于 \(CD\)，所以 \(CD\perp\) 平面 \(BEF\). 又 \(CD\subset\) 平面 \(PCD\)，故平面 \(BEF\perp\) 平面 \(PCD\).
\end{enumerate}
\end{solution}


\begin{problem}
已知函数 \( f(x) = x^2 + x \sin x + \cos x \).
\begin{enumerate}
\item 若曲线 \( y = f(x) \) 在点 \((a, f(a))\) 处与直线 \( y = b \) 相切，求 \( a \) 与 \( b \) 的值；
\item 若曲线 \( y = f(x) \) 与直线 \( y = b \) 有两个不同交点，求 \( b \) 的取值范围.
\end{enumerate}
\end{problem}

\begin{solution}
\begin{enumerate}
\item \(f'(x)=2x+\sin x+x\cos x-\sin x=x(2+\cos x)\). 因为 \(-1\leqslant\cos x\leqslant1\)，所以 \(2+\cos x>0\). 直线 \(y=b\) 是水平直线，曲线在 \((a,f(a))\) 处与其相切，故切点满足 \(f'(a)=0\) 且 \(f(a)=b\). 由 \(a(2+\cos a)=0\) 及 \(2+\cos a>0\)，得 \(a=0\)，进而 \(b=f(0)=1\).
\item 由 \(f'(x)=x(2+\cos x)\) 及 \(2+\cos x>0\)，知当 \(x<0\) 时 \(f'(x)<0\)，当 \(x>0\) 时 \(f'(x)>0\). 因此 \(f(x)\) 在 \((-\infty,0)\) 上单调递减，在 \((0,+\infty)\) 上单调递增，且 \(f(0)=1\) 是最小值. 又因为 \(f(x)\geqslant x^2-|x|-1\)，当 \(x\to\pm\infty\) 时 \(f(x)\to+\infty\). 所以当 \(b>1\) 时，直线 \(y=b\) 分别与左右两个单调区间的图象各有一个交点；当 \(b=1\) 时只有 \(x=0\) 一个交点；当 \(b<1\) 时没有交点. 故曲线与直线有两个不同交点时，\(b\in(1,+\infty)\).
\end{enumerate}
\end{solution}

\begin{problem}
直线 \( y = kx + m  \) （\(m \neq 0\)）与椭圆 \( W: \frac{x^2}{4} + y^2 = 1 \) 相交于 \(A\)，\(C\) 两点，\(O\) 是坐标原点.
\begin{enumerate}
\item 当点 \(B\) 的坐标为 \((0,1)\)，且四边形 \(OABC\) 为菱形时，求 \(AC\) 的长；
\item 当点 \(B\) 在 \(W\) 上且不是 \(W\) 的顶点时，证明：四边形 \(OABC\) 不可能为菱形.
\end{enumerate}
\end{problem}

\begin{solution}
\begin{enumerate}
\item 菱形 \(OABC\) 的两条对角线互相垂直平分，故 \(OB\perp AC\)，且 \(AC\) 的中点为 \(OB\) 的中点. 由 \(O=(0,0)\)，\(B=(0,1)\)，知 \(OB\) 在 \(y\) 轴上，所以 \(AC\) 是水平弦，且其中点为 \(\left(0,\frac{1}{2}\right)\). 因而点 \(A\)，\(C\) 的纵坐标均为 \(\frac{1}{2}\). 代入椭圆方程得 \(\frac{x^2}{4}+\frac{1}{4}=1\)，所以 \(x=\pm\sqrt{3}\). 因此 \(A\)，\(C\) 的坐标为 \(\left(\pm\sqrt{3},\frac{1}{2}\right)\)，从而 \(AC=2\sqrt{3}\).
\item 假设四边形 \(OABC\) 为菱形. 因为 \(OA=OC\)，设 \(OA=OC=r\)，则 \(A\)，\(C\) 同时满足圆方程 \(x^2+y^2=r^2\) 和椭圆方程 \(\frac{x^2}{4}+y^2=1\). 两式相减，得 \(\frac{3}{4}x^2=r^2-1\)，所以 \(x_A^2=x_C^2\)，即 \(x_A=x_C\) 或 \(x_A=-x_C\).

若 \(x_A=x_C\)，由于 \(A\neq C\)，直线 \(AC\) 为竖直直线. 菱形的对角线互相垂直，所以 \(OB\) 为水平直线，进而 \(y_B=0\). 又 \(B\) 在椭圆 \(W\) 上，故 \(\frac{x_B^2}{4}=1\)，于是 \(B=(2,0)\) 或 \(B=(-2,0)\)，这两个点都是椭圆的顶点.

若 \(x_A=-x_C\)，由菱形对角线互相平分，点 \(B\) 是线段 \(AC\) 中点的两倍位置，即 \(\frac{B}{2}\) 为 \(AC\) 的中点. 因此 \(x_B=x_A+x_C=0\). 再由 \(B\in W\)，得 \(y_B^2=1\)，于是 \(B=(0,1)\) 或 \(B=(0,-1)\)，同样是椭圆的顶点.

两种情形都与点 \(B\) 不是 \(W\) 的顶点矛盾，所以四边形 \(OABC\) 不可能为菱形.
\end{enumerate}
\end{solution}

\begin{problem}
给定数列 \(a_1\)，\(a_2\)，\(\cdots\)，\(a_n\)，对 \(i = 1\)，\(2\)，\(\cdots\)，\(n-1\)，该数列前 \(i\) 项的最大值记为 \(A_i\)，后 \(n-i\) 项 \(a_{i+1}\)，\(a_{i+2}\)，\(\cdots\)，\(a_n\) 的最小值记为 \(B_i\)，\(d_i = A_i - B_i\).
\begin{enumerate}
\item 设数列 \(\{a_n\}\) 为 3， 4， 7， 1，写出 \(d_1\)，\(d_2\)，\(d_3\) 的值；
\item 设 \(a_1\)，\(a_2\)，\(\cdots\)，\(a_n\ \) （\(n \geqslant 4\)）是公比大于 1 的等比数列，且 \(a_1 > 0\). 证明： \(d_1\)，\(d_2\)，\(\cdots\)，\(d_{n-1}\) 是等比数列；
\item 设 \(d_1\)，\(d_2\)，\(\cdots\)，\(d_{n-1}\) 是公差大于 0 的等差数列，且 \(d_1 > 0\). 证明： \(a_1\)，\(a_2\)，\(\cdots\)，\(a_{n-1}\) 是等差数列.
\end{enumerate}
\end{problem}

\begin{solution}
\begin{enumerate}
\item 当 \(i=1\) 时，\(A_1=3\)，\(B_1=\min\{4,7,1\}=1\)，所以 \(d_1=3-1=2\). 当 \(i=2\) 时，\(A_2=\max\{3,4\}=4\)，\(B_2=\min\{7,1\}=1\)，所以 \(d_2=4-1=3\). 当 \(i=3\) 时，\(A_3=\max\{3,4,7\}=7\)，\(B_3=1\)，所以 \(d_3=7-1=6\).
\item 设等比数列的公比为 \(q\). 因为 \(a_1>0\) 且 \(q>1\)，所以 \(a_i=a_1q^{i-1}\) 且 \(a_1<a_2<\cdots<a_n\). 因而对 \(i=1\)，\(2\)，\(\cdots\)，\(n-1\)，有 \(A_i=a_i\)，\(B_i=a_{i+1}\)，从而 \(d_i=a_i-a_{i+1}=a_1q^{i-1}(1-q)\). 于是对 \(i=1\)，\(2\)，\(\cdots\)，\(n-2\)，\(d_i\neq0\)，且 \(\frac{d_{i+1}}{d_i}=q\) 为常数，所以 \(d_1\)，\(d_2\)，\(\cdots\)，\(d_{n-1}\) 是等比数列.
\item 设 \(d_1\)，\(d_2\)，\(\cdots\)，\(d_{n-1}\) 的公差为 \(d\)，则 \(d>0\)，从而 \(d_1<d_2<\cdots<d_{n-1}\). 若 \(n=2\)，结论显然. 以下设 \(n\geqslant3\).

先证 \(a_1\)，\(a_2\)，\(\cdots\)，\(a_{n-1}\) 递增. 若不然，设 \(k\) 是满足 \(2\leqslant k\leqslant n-1\) 且 \(a_k\leqslant a_{k-1}\) 的最小下标. 由 \(a_1<a_2<\cdots<a_{k-1}\)，得 \(A_k=A_{k-1}=a_{k-1}\). 又因为 \(B_{k-1}=\min\{a_k,\cdots,a_n\}\leqslant B_k=\min\{a_{k+1},\cdots,a_n\}\)，所以 \(d_{k-1}=A_{k-1}-B_{k-1}\geqslant A_k-B_k=d_k\)，这与 \(d_{k-1}<d_k\) 矛盾. 因此 \(a_1<a_2<\cdots<a_{n-1}\).

由 \(d_1>0\) 及 \(A_1=a_1\)，得 \(B_1=A_1-d_1=a_1-d_1<a_1\). 但 \(a_2\)，\(a_3\)，\(\cdots\)，\(a_{n-1}\) 都大于 \(a_1\)，而 \(B_1\) 是 \(a_2\)，\(a_3\)，\(\cdots\)，\(a_n\) 的最小值，所以必有 \(a_n<a_1\). 因而对任意 \(i=1\)，\(2\)，\(\cdots\)，\(n-1\)，有 \(A_i=a_i\)，\(B_i=a_n\)，从而 \(a_i=a_n+d_i\). 由于 \(d_i\) 是等差数列，故 \(a_1\)，\(a_2\)，\(\cdots\)，\(a_{n-1}\) 也是等差数列，且公差为 \(d\).
\end{enumerate}
\end{solution}
